% Glyphs that Latin Modern lacks, kept so the document can stay in the
% handouts' typeface instead of switching to a Unicode-wide font.
%
% Two different problems. The Greek letters and operators arrive as literal
% Unicode inside ordinary prose, because the questions were pasted from a
% chat window; they are rendered as the maths they mean. The accented "u" is
% the opposite case: it is fine in text but reaches maths mode inside Italian
% words such as "piu`", where the maths italic font has no accents, so it is
% pushed back into a text box when that happens.
\usepackage{newunicodechar}
\newunicodechar{α}{\ensuremath{\alpha}}
\newunicodechar{θ}{\ensuremath{\theta}}
\newunicodechar{∥}{\ensuremath{\|}}
\newunicodechar{≠}{\ensuremath{\neq}}
\newunicodechar{ù}{\ifmmode\text{\`u}\else\`u\fi}

% Tighten the space around displayed equations. The transcript puts almost every
% formula on its own display line - 2,427 of them, many a single symbol - and
% LaTeX's default skips around each turn that habit into roughly a hundred extra
% pages. This changes spacing only: no equation is reflowed and no word is cut.
\usepackage{amsmath}
\setlength{\abovedisplayskip}{3pt plus 2pt minus 2pt}
\setlength{\belowdisplayskip}{3pt plus 2pt minus 2pt}
\setlength{\abovedisplayshortskip}{1pt plus 1pt}
\setlength{\belowdisplayshortskip}{2pt plus 1pt}
\setlength{\parskip}{2pt plus 1pt}
